\subsection{Quelltextauszug zu Datenschutz und Sicherheit}
\label{app:CodeAuthentifizierung}

Die folgenden Auszüge zeigen die sicherheitsrelevanten Kernfunktionen der
Anwendung. Sie sind nicht als vollständiger Quelltextabdruck gedacht, sondern
folgen dem Ablauf eines geschützten Pflegevorgangs: Anmeldung, Prüfung der
Berechtigung, Änderung von Daten, Protokollierung und optionales
Zurücksetzen.

\subsubsection*{Anmeldung und Session}
Der Einstieg in die Anwendung erfolgt über die Login-Funktion. Die E-Mail-Adresse
wird normalisiert, danach werden Rate-Limit, Kontosperre und Passwortprüfung
durchlaufen. Erst bei erfolgreicher Prüfung wird eine neue Session gestartet.

\lstinputlisting[
  language=Python,
  firstline=25,
  lastline=91,
  firstnumber=25,
  caption={Authentifizierung und Start der Benutzersession}
]{PROJEKT/app/services/authentication.py}
\clearpage

\subsubsection*{Passwortschutz}
Passwörter werden nicht im Klartext gespeichert. Beim Anlegen oder Ändern eines
Benutzers wird das Passwort gehasht. Zusätzlich gelten Mindestanforderungen an
die Passwortstärke.

\lstinputlisting[
  language=Python,
  firstline=1,
  lastline=27,
  firstnumber=1,
  caption={Hashing und Prüfung der Passwortstärke}
]{PROJEKT/app/services/passwords.py}
\clearpage

\subsubsection*{Schutz gegen wiederholte Fehlanmeldungen}
Die Login-Sicherheit arbeitet auf zwei Ebenen. Wiederholte Fehlversuche von
einer IP-Adresse werden zeitlich begrenzt gezählt. Zusätzlich wird je Konto
gespeichert, ob eine Sperre aktiv ist. Jeder Login-Versuch wird protokolliert.

\lstinputlisting[
  language=Python,
  firstline=19,
  lastline=89,
  firstnumber=19,
  caption={Rate-Limit, Kontosperre und Login-Audit}
]{PROJEKT/app/services/login_security.py}
\clearpage

\subsubsection*{Rollen und Mandantentrennung}
Die Mandantentrennung wird nicht nur in der Oberfläche abgebildet. Die
Service-Funktionen filtern Datensätze anhand der Rolle und der zugeordneten
Firma. Benutzer ohne Root-Rolle sehen dadurch nur Daten aus ihrer
Firma.

\lstinputlisting[
  language=Python,
  firstline=8,
  lastline=69,
  firstnumber=8,
  caption={Rollenprüfung und Filterung sichtbarer Daten}
]{PROJEKT/app/services/permissions.py}
\clearpage

\subsubsection*{Schutz vor firmenfremden Änderungen}
Vor jeder Änderung wird geprüft, ob der angemeldete Benutzer den betroffenen
Datensatz bearbeiten darf. Root darf firmenübergreifend arbeiten,
andere Benutzer bleiben auf ihre Firma begrenzt. Root-Benutzer werden
zusätzlich vor versehentlicher Bearbeitung durch andere Rollen geschützt.

\lstinputlisting[
  language=Python,
  firstline=78,
  lastline=111,
  firstnumber=78,
  caption={Prüfung bearbeitbarer Firmen und Benutzer}
]{PROJEKT/app/services/permissions.py}
\clearpage

\subsubsection*{Auditierung und Schutz sensibler Daten}
Fachliche Änderungen werden im Audit-Log gespeichert. Dabei werden sensible
Felder wie Passwörter vor dem Speichern maskiert. Dadurch bleibt die Änderung
nachvollziehbar, ohne Passwortwerte im Audit-Log offenzulegen.

\lstinputlisting[
  language=Python,
  firstline=9,
  lastline=49,
  firstnumber=9,
  caption={Maskierung sensibler Felder im Audit-Log}
]{PROJEKT/app/services/audit.py}

\lstinputlisting[
  language=Python,
  firstline=121,
  lastline=145,
  firstnumber=121,
  caption={Erzeugen eines Audit-Eintrags}
]{PROJEKT/app/services/audit.py}
\clearpage

\subsubsection*{Audit-Einträge für Änderungen}
Für Erstellen, Ändern und Löschen werden jeweils passende Audit-Informationen
abgelegt. Bei Änderungen wird sowohl der vorherige als auch der neue Zustand
gespeichert. Dadurch kann später nachvollzogen werden, was konkret geändert
wurde.

\lstinputlisting[
  language=Python,
  firstline=148,
  lastline=213,
  firstnumber=148,
  caption={Audit-Einträge für Erstellen, Ändern und Löschen}
]{PROJEKT/app/services/audit.py}
\clearpage

\subsubsection*{Kontrolliertes Zurücksetzen}
Für ausgewählte Audit-Einträge kann der vorherige Zustand wiederhergestellt
werden. Auch dieser Vorgang erzeugt wieder einen neuen Audit-Eintrag, damit
das Zurücksetzen selbst nachvollziehbar bleibt.

\lstinputlisting[
  language=Python,
  firstline=286,
  lastline=365,
  firstnumber=286,
  caption={Rollback auf Basis eines Audit-Eintrags}
]{PROJEKT/app/services/audit.py}
\clearpage

\subsubsection*{API-Beispiel}
Die API nutzt dieselben Berechtigungsprüfungen wie die Weboberfläche. Im
Beispiel dürfen nur Manager eine Mail-Adresse anlegen. Beim Lesen werden die
sichtbaren Mail-Adressen über den aktuell angemeldeten Benutzer gefiltert.

\lstinputlisting[
  language=Python,
  firstline=21,
  lastline=54,
  firstnumber=21,
  caption={API-Auszug für Mail-Adressen mit Rollenprüfung}
]{PROJEKT/app/web/api/mail_addresses.py}
